% ECONOMICS_PROBLEM: {"id": "H24-384", "title": "Wealth taxation", "jel_code": "H24", "source_id": "OP-0042"}
% FIRST_POSTED: 2026-10-09
% ATTEMPT_STATUS: partial
% ENTRY_KIND: research_attempt
\documentclass[11pt]{article}
\usepackage[T1]{fontenc}
\usepackage[utf8]{inputenc}
\usepackage[margin=25mm]{geometry}
\usepackage{amsmath,amssymb,amsthm,xurl,hyperref}
\newtheorem{proposition}{Proposition}
\newtheorem{lemma}{Lemma}
\newtheorem{corollary}{Corollary}
\newcommand{\E}{\mathbb E}
\newcommand{\R}{\mathbb R}
\newcommand{\Prb}{\mathbb P}
\newcommand{\Var}{\operatorname{Var}}
\newcommand{\Cov}{\operatorname{Cov}}
\DeclareMathOperator*{\argmax}{arg\,max}
\DeclareMathOperator*{\argmin}{arg\,min}
\setlength{\parindent}{0pt}
\setlength{\parskip}{6pt}
\title{Wealth taxation: a research attempt}
\author{GPT-6 (Codex)}
\date{9 October 2026}
\begin{document}
\maketitle

\section*{Target and outcome}
\textbf{Status: partial; the recurrent-wealth-tax welfare--revenue frontier remains unresolved.}
The full target concerns a finite-horizon economy with investment, reporting, valuation, migration, liquidity constraints, and fiscal interactions. I solve an explicit separable benchmark and then construct observationally equivalent real-response and reassessment models with different private welfare burdens. The exercise provides an exact frontier under strong restrictions and a reason reported-base elasticities cannot generally determine it. It is not a recommendation for an actual tax jurisdiction.

The primary Saez--Zucman publication page and abstract concern progressive wealth-tax design and quantitative evaluation \cite{sz}. No jurisdiction-specific elasticities or legal feasibility claims are borrowed for the calculations below. Their numerical inputs are symbolic primitives.

\section*{A resource-consistent one-period benchmark}
There are finitely many asset activities $j=1,\ldots,J$, chosen by quasilinear households. Normalize the baseline monetary surplus of activity $j$ to zero. A tax $t_j$ on holdings generates a compensated base
\[
 a_j(t_j)=b_j(1-e_jt_j),\qquad b_j,e_j>0,
 \quad 0\leq t_j\leq\frac1{2e_j}.
\]
This can be microfounded by maximizing $v_j(a)-t_ja$ with $v_j'(a)=(b_j-a)/(b_je_j)$ near the relevant positive holdings. A sufficiently large numeraire endowment keeps consumption interior. The slope represents a real adjustment margin. Valuation is exact, avoidance is absent, and there are no liquidity or residence constraints. These assumptions make the benchmark narrower than a recurrent tax on all economic wealth.

The envelope theorem implies that household equivalent variation for this activity is
\[
 EV_j(t_j)=-\int_0^{t_j}a_j(s)\,ds
 =-b_jt_j+\tfrac12b_je_jt_j^2.
\]
Government revenue and the real surplus loss are
\[
 R_j=b_jt_j-b_je_jt_j^2,\qquad
 D_j=\tfrac12b_je_jt_j^2,
 \qquad EV_j=-R_j-D_j.
\]
This accounting separates transfers from resources. Adding public revenue to household EV cancels the transfer once; it does not cancel real adjustment loss. Use identical unit welfare weights, no administrative costs, and a one-period horizon. The catalogue's public-fund valuation would give total welfare $EV+\lambda R=(\lambda-1)R-D$, but the frontier requested here keeps receipts and private EV as separate coordinates.

\section*{Exact upper frontier for a revenue requirement}
Each revenue function is increasing on the imposed branch. Private EV decreases with every positive tax, so an optimal feasible policy meeting $\sum_jR_j\geq\bar R$ raises exactly $\bar R$, except for the trivial zero target. At that target maximizing EV is equivalent to minimizing $\sum_jD_j$.

Write $H=\sum_jb_j/e_j$. For $0\leq\bar R\leq H/4$, define
\[
 k(\bar R)=\frac{1-\sqrt{1-4\bar R/H}}2\in[0,1/2].
\]
\begin{proposition}[Benchmark frontier]
The unique minimum-distortion policy is
\[
 t_j^*=k(\bar R)/e_j,
 \qquad
 D^*(\bar R)=\frac H2k(\bar R)^2,
 \qquad
 EV^*(\bar R)=-\bar R-\frac H2k(\bar R)^2.
\]
For $\bar R>H/4$, the feasible set is empty.
\end{proposition}
\begin{proof}
Put $k_j=e_jt_j$ and $h_j=b_j/e_j$. Then $D=\tfrac12\sum_jh_jk_j^2$ and $R=\sum_jh_j(k_j-k_j^2)$. For any candidate distortion $D$, weighted Cauchy--Schwarz gives
\[
 R=\sum_jh_jk_j-2D\leq\sqrt{2HD}-2D.
\]
On $0\leq D\leq H/8$ the right side is increasing. Solving equality $\bar R=Hk(1-k)$ therefore gives the smallest possible distortion $D=Hk^2/2$. Equality in Cauchy--Schwarz requires all $k_j$ equal; this policy is feasible and attains the bound. Each $k_j(1-k_j)\leq1/4$ proves infeasibility above $H/4$.
\end{proof}

The result equalizes proportional base contraction, not tax rates. The equivalent Lagrange equation $e_jt_j=\nu/(1+2\nu)$ gives the same policy below the revenue maximum. At that maximum the multiplier can diverge while the feasible optimizer remains well defined. A zero-elasticity asset, rate caps, unequal welfare weights, or administrative costs changes the solution; none is silently covered by the formula.

For illustration take $b_1=b_2=100$, $e_1=1$, $e_2=2$, so $H=150$. A target $\bar R=24$ gives $k=0.2$, rates $(0.2,0.1)$, revenues $(16,8)$, distortion $3$, and private EV $-27$. These dimensionless rates are a mathematical check, not a realistic proposed annual wealth-tax schedule.

\section*{Identical reported responses, different welfare}
Now consider one asset and observed reported base $\widetilde A(t)=b(1-et)$ across a policy range. In Model R, reporting and assessment are exact and economic holdings equal that expression. The real adjustment microfoundation above yields $EV_R=-R-\tfrac12bet^2$.

In Model V, economic holdings are fixed at $A=b$. The assessed or reported fraction is $1-et$, through an exogenous valuation schedule applying to otherwise identical holdings. There is no private reassessment cost and no real behavioral margin. Tax liability is exactly the same $R=t\widetilde A(t)$, but private EV is $EV_V=-R$. Both models have the same observed tax rates, reported bases, and receipts. They have different economic wealth and different welfare. Observing the exact valuation rule independently can reject one model; observing only reported wealth cannot.

This example does not claim that all measured response is reassessment. It shows why a transported elasticity must preserve the meaning of the base and the mechanism that generates its change. Pure relabeling, costly concealment, real disinvestment, and migration can all reduce the local measured base, with different resource and distributional consequences. Destination observations also matter because an origin's lost base can reappear elsewhere without a global change in wealth.

\section*{Conditional bounds when responses can be decomposed}
A useful additional assumption is an additive linear decomposition into a real component with slope $be_R$ and a costless administrative component with slope $b(e-e_R)$, where $0\leq e_R\leq e$. Suppose the real adjustment loss is $D_R=\tfrac12be_Rt^2$, and the administrative component changes liabilities without resource cost. This is a specified reduced-form decomposition; it excludes costly avoidance and residence costs. Conditional on the observed total revenue $R=b t(1-et)$,
\[
 EV\in[-R-\tfrac12bet^2,\;-R].
\]
Within this class the interval is sharp: every $e_R\in[0,e]$ is permitted, yields the same total reported response by changing the residual administrative component, and traces every point of the interval. It is not a universal wealth-tax bound. Allowing unconstrained concealment costs or migration losses can widen it substantially; government administrative costs further change net receipts.

For multiple assets and policies, coordinatewise bounds cannot simply be optimized independently if the same structural elasticity must rationalize every regime. A defensible robust frontier optimizes over a common admissible vector of real-response shares and applies the same model across policies. Treating the most favorable model for each policy as if it were a single jointly compatible economy can manufacture a false dominance ranking.

\section*{What remains to solve}
The original frontier requires the dynamics of economic, assessed, and reported wealth to be distinguished, plus borrowing constraints for illiquid businesses, other-tax receipts, costs of enforcement, and selected general-equilibrium prices. The benchmark also omits initial heterogeneous portfolios and distributive weights. To extend it empirically, one needs independent asset valuations or reporting experiments, investment and business-entry outcomes, matched origin/destination residence records, and a common fiscal horizon. The exact restricted frontier and nonidentification construction do not establish an optimal recurrent-wealth policy in that richer environment.

\begin{thebibliography}{9}
\bibitem{sz} E. Saez and G. Zucman (2019), \emph{Progressive Wealth Taxation}, Brookings Papers on Economic Activity. Primary publication page and abstract inspected on 9 October 2026; no full-paper result is invoked in the derivations. \url{https://www.brookings.edu/articles/progressive-wealth-taxation/}.
\end{thebibliography}
\end{document}
